\documentclass[11pt]{article}
\usepackage[margin=0.98in]{geometry}
\usepackage{amsmath,amssymb,hyperref}
\hypersetup{colorlinks=true,urlcolor=blue}
\title{A finite-quotient replacement for the\\ explicit unitary class count}
\author{Proof-shortening proposal for ProveIt}
\date{28 November 2026}
\newcommand{\OO}{\mathcal O}
\newcommand{\CC}{\mathbb C}
\newcommand{\QQ}{\mathbb Q}
\begin{document}
\maketitle
\noindent\textbf{Result.}
The proof of \texttt{New.Space.card\_stdClasses\_le} can avoid bounded
strong approximation and the auxiliary matrix $GX$. Keep the short basis
and its Gram matrix $G$. At each $G$, count residues of the adelic inverse
in the finite quotient $G^{-1}R^r/R^r$. The current explicit constant and
the downstream theorem statement are unchanged. This is a written
replacement, not a Lean-checked patch or a measurement of lines saved.

\section*{1. The replacement invariant}
Use the notation and hypotheses of the
\href{https://github.com/VladimirReshetnikov/ProveIt/blob/2c38f343aa14c8ce85ab1c7a203774c7d05a9780/Analysis/GeneralizedRiemannHypothesis/Lean/grh/src/New/UpperBound/Size/ClassCount.lean#L116}{pinned class-count theorem}.
Let $E$ be the CM field, $R=\widehat{\OO}_E$ its ring of integral finite
adeles, and $r$ the rank. For each standard double coset $g$ choose the
same short basis $Y$ furnished by the existing proof. Over the finite adeles
put
\[
 A=g_{\rm out}^{-1}Y,\qquad G=Y^\ast Y=A^\ast A .
\]
Then $A$ has entries in $R$, $G$ is invertible over the adeles, and $G$ has
entries in $\OO_E$ by the existing integrality lemma. Define
\[
 L_G=\{v\in\mathbb A_{E,f}^{r}:Gv\in R^r\}=G^{-1}R^r .
\]
It contains $R^r$. Every column of $A^{-1}$ belongs to $L_G$ because
$GA^{-1}=A^\ast$. Record $G$ and the $r$ ordered column classes of
$A^{-1}$ in $L_G/R^r$. Equivalently record the columns of $A^\ast$ in
$R^r/GR^r$. A class representative and a short basis may be chosen once;
independence from all choices is not needed for an injection.

\paragraph{Collision.}
For equal data at $g,g'$ the difference $A^{-1}-A'^{-1}$ is integral.
Since $A'$ is integral,
\[
 A'A^{-1}=I+A'(A^{-1}-A'^{-1})
\]
is integral. With the equality of rational Gram matrices, this is exactly
the hypothesis of\par\noindent \texttt{stdClasses\_eq\_of\_gram\_eq}. Thus
$g=g'$. No approximation or lifting of arbitrary residue tuples is needed.

\section*{2. Cardinality and the explicit constant}
Multiplication by $G$ identifies $L_G/R^r$ with $R^r/GR^r$.
At each finite place the local integer ring is a DVR; integral row and
column operations diagonalize the full-rank matrix $G$ there. Its local
cokernel has size $q_{\mathrm{p}}^{v_{\mathrm{p}}(\det G)}$, where
$q_{\mathrm{p}}$ is the residue-field size. Almost all factors are $1$.
Taking the product gives
\[
 |L_G/R^r|=|N_{E/\QQ}(\det G)|.
\]
In Lean the local diagonalization need not be added: see the route below.

Let $f$ be the number of infinite places and $c=6N(fr)^2$. For $r>0$ the
short-basis bound gives $\lvert\phi(Y_{ij})\rvert\le c$ at every
embedding $\phi$. Set $W=rc^2$. Every Gram entry is bounded by $W\le c^3$.
The existing integer box estimate gives at most $(1+2c^3)^{2fr^2}$
Gram matrices. The Leibniz formula yields
\[
 |\phi(\det G)|\le r!W^r\le(rW)^r\le c^{3r}
\]
since $r^2\le c$. The conjugate pairs of complex embeddings give
$|N_{E/\QQ}(\det G)|^r\le c^{6fr^2}$ possible column tuples per Gram.
Multiply the bounds, using $c^3\le 1+2c^3$, to get exactly the current
$(1+2c^3)^{4fr^2}$ upper bound. Rank zero has one class.

\paragraph{Optional narrower count.}
Hadamard on $\phi(Y)$ gives $|\phi(\det G)|=|\det\phi(Y)|^2\le W^r$.
Keeping $W$ in the box estimate gives
$(1+2W)^{2fr^2}W^{2fr^2}$ for positive rank.
This improvement is unnecessary to retain the existing interface.

\section*{3. Lean integration and verification boundary}
Retain the short-basis choice, the Gram-entry integrality lemma, the
integer box estimate, and the final equality lemma. Replace the approximation
and second bounded-matrix construction as follows.

\smallskip\noindent\textbf{Quotient.}
Define the fiber for fixed $G$ as the quotient attached to
\texttt{preimageLattice} by \texttt{intVectors} (i.e.\ $L_G/R^r$).
For the column of $A^{-1}$, membership follows directly from $GA^{-1}=A^\ast$.
The quotient equivalence relation gives an integral column difference.
Use a dependent sum over bounded integral Gram matrices, or count
each fiber of the Gram map separately.

\smallskip\noindent\textbf{Index.}
Identify the preimage lattice with the column lattice of $G^{-1}$
by applying $G$ and $G^{-1}$. Then
\href{https://github.com/VladimirReshetnikov/ProveIt/blob/2c38f343aa14c8ce85ab1c7a203774c7d05a9780/Analysis/GeneralizedRiemannHypothesis/Lean/grh/src/Lib/General/NumberTheory/LatticeCovolumes.lean#L197}{\texttt{submoduleIndex\_intVectors\_mul\_genIndex\_det}}
computes the relative index using $\det(G^{-1})$.
Alternatively use \texttt{submoduleIndex\_colLattice\_div} on $G$.
Combine multiplicativity of the generalized index with
\href{https://github.com/VladimirReshetnikov/ProveIt/blob/2c38f343aa14c8ce85ab1c7a203774c7d05a9780/Analysis/GeneralizedRiemannHypothesis/Lean/grh/src/Lib/Adelic/Adeles/IdeleModulus.lean#L399}{\texttt{New.genIndex\_algebraMap}}.
For an implementation, establish the nonzero relative index and identify
it with the relevant quotient's \texttt{Nat.card} before counting tuples.
Prove the determinant bound with the ordinary finite sum over permutations;
a Hadamard API is not required.

\smallskip\noindent\textbf{What disappears.}
The bounded principal representative $X$, both bounds on $X$ and $GX$,
and the corresponding integrality and collision calculations are
local to the current argument. No claim is made that all imports used
for them disappear from the wider development. Lean elaboration and net
line savings have not been measured.

\smallskip\noindent\textbf{Deduplication.}
The accompanying audit pins the public source and proposal files.
I found no proposal for this invariant in that baseline. The included
exact integer-matrix replay tests the two congruence descriptions and the
collision calculation, including noncollision cases. It is a check of the
algebra, not a test of the finite-adelic Lean proof.
\end{document}
